% D.2.48 working proof.  Do not register until the full peer audit.
% The refined-class argument below does not depend on a global covering.
\paragraph{Theorem D.2.48 (Proved, pending independent audit: the weighted Engel refined class).}
Take the weight-$d$ Engel construction of D.2.34, for any $d>0$,
with both additive sites replaced by good logarithmic caps.  Retain
all the actual marked data
\[
 m_0=3k,\qquad m_1=4\ell,\qquad
 v_i=r_i\delta+n_i u_i,\qquad
 P=m_1n_0+m_0n_1=\pm1 .
 \tag{D.WI1}
\]
The manifold is simply connected.  Its complex first and second
Chern classes, spin characteristic class, and refined class satisfy
\[
 c_1(X)=c_2(X)=\lambda(X)=0,\qquad P_{\mathrm{ref}}(X)=0.
 \tag{D.WI2}
\]
Here $P_{\mathrm{ref}}$ denotes Zhubr's class, to distinguish it from
the Seifert determinant $P$ in (D.WI1).

The classes $x_i=\operatorname{PD}(D_i)$ of the $d$ cusp components,
indexed cyclically, generate $H^2(X;\mathbb Z)$ with sole relation
$\sum_i x_i=0$.  The complex ruling $f=H-E_1$ has image
$s=[a\wedge b]$, a generator of $\operatorname{Tor}H_2(X)=\mathbb Z/d$.
Put $\kappa=\operatorname{PD}(f)$.  With repetitions retained when
$d=2$, the complete degree-two product is
\[
 x_i x_j=
 \bigl(\mathbf1_{j=i+1}+\mathbf1_{j=i-1}-2\mathbf1_{j=i}\bigr)\kappa.
 \tag{D.WI3}
\]
For $d=1$ the sole component class is zero.  For $d=2$ the two
off-diagonal incidences both occur, and every product in (D.WI3)
is zero since $2\kappa=0$.  The integral cubic is zero for every $d$.
These assertions do not by themselves specify the complete
fundamental-class invariant at primes two and three.

\begin{proof}
\emph{Primitive component classes.}
Write $E$ for the complement of the interior of the cusp tube.
The actual finite-cap presentation before adjoining its cusp meridian is
\[
 H_1(E)=
 \langle\gamma,x_0',x_1'\mid m_0x_0'=n_0\gamma,\ 
                                   m_1x_1'=n_1\gamma\rangle .
\]
The determinant condition implies
$\gcd(m_0,m_1)=\gcd(m_i,n_i)=1$ and $H_1(E)=\mathbb Z$.
The cusp meridian is $x_0'+x_1'$ up to the common orientation sign.
It is primitive because adjoining its relation gives $H_1(X)=0$.
Angular differences of adjacent toric meridians are killed in
$H_1(E)$.  Thus the relative Thom map is
\[
 H^1(E)\longrightarrow H^2(X,E)=\mathbb Z^d,\qquad
                   1\longmapsto\pm(1,\ldots,1).
 \tag{D.WI4}
\]
The cokernel of the resulting injection
$\mathbb Z^d/\mathbb Z(1,\ldots,1)\hookrightarrow H^2(X)$
embeds in $H^2(E)$, which is torsion-free by the integral universal
coefficient theorem.  It has rank zero by D.2.34, so it vanishes.
In particular the entire fibre class is
$f_X^*[\mathrm{point}]=\sum_i x_i=0$.
For a reduced multiple fibre $D$ one has
$m[D]=f_X^*[\mathrm{point}]=0$; the group $H^2(X)$ is free,
so $[D]=0$ as well.

The common ruling in the $d$-period cusp is the specialization
of $a\wedge b$, with its complex orientation.  This is an integral
unit-period annulus calculation: that torus remains in one gamma
component and limits to the two adjacent external curves
$E_2+E_3=H-E_1$.  The three support functions
$r(r-1)/2$, $rs$, and $s(s-1)/2$ evaluate on that ruling as
$(-1,0,0)$, exactly their actual multiplier windings on
$a\wedge b$; its dual-torus coordinate is zero.
There is no division by $d$ or saturation at this step.
The primitive functional of D.2.34 evaluates it as
$\xi_d(a\wedge b)=1$.  The matrix (D.WE3) shows that $s$ has
exact order $d$ in the simply connected sector.  The actual
external intersections and the principal sum of components
therefore give (D.WI3).  The adjacent-pair formula counts both
edges when they have the same pair of labels.  Since every
product is torsion, its further product with $H^2(X)$ vanishes
in $H^6(X)=\mathbb Z$.

\emph{The differential and its literal critical cycle.}
The normalization immersion $j_i:dP_6\longrightarrow X$ has
normal line $K_{dP_6}$.  Indeed the principal local central
divisor is the product of the normal coordinates to all its
branches; on a fixed normalized branch its normal transition
is the inverse of the six toric boundary transitions.  This
includes both conductor branches of each self-identification.
Consequently
\[
 j_i^*c_1(TX)=c_1(TdP_6)+c_1(K_{dP_6})=0,\qquad
 \langle j_i^*c_2(TX),[dP_6]\rangle=6-6=0 .
 \tag{D.WI5}
\]
Restriction of $H^2(X)$ to the normalized components is injective.
For a component combination with coefficients $u_i$, its
restriction to component $i$ is
$\sum_r(u_{i+r}-u_i)C_r$.  Pairing with the anticanonical ample
class, which has degree one on each of the six boundary curves,
gives the component-graph Laplacian.  The graph is connected, so
zero restrictions force the $u_i$ to be constant.  That is
precisely the already removed principal relation.  Equation
(D.WI5) now gives $c_1(TX)=0$ integrally.

Remove the divisorial multiple-fibre factors from $df_X$.  Its
corrected target is the line
\[
 L=f_X^*T\mathbb P^1\otimes
                      \mathcal O\bigl(-\sum (m_i-1)D_i^{\rm fin}\bigr).
\]
It is topologically trivial by (D.WI4).  The abstract rank-two
kernel $V$ is locally free, including at $xy$ and $xyz$ charts,
and has the nowhere-zero abstract section given by $\delta$.
At a node the kernel columns are $(x,-y,0)$ and $(0,0,1)$;
at a triple point they are $(x,-y,0)$ and $(x,0,-z)$.
The constant scalar-circle weight vector is a nonzero constant
combination in these abstract bases.  At the good finite caps
the corrected differential is a submersion and $\delta$ remains
a nonzero tangent vector.  Since
$c_1(V)=c_1(TX)-c_1(L)=0$, its quotient by this section is a
trivial complex line.  Thus $V$ and $L$ admit full smooth frames.
They give the reference frame $\rho$ on the complement of the
literal critical curve cycle $C$.

The local Bott clutch at $xy$ has coefficient one:
its two normal columns are
$\bigl(\begin{smallmatrix}x&\overline y\\-y&\overline x\end{smallmatrix}\bigr)$.
The $xyz$ Hilbert--Burch minors give multiplicity one on each
of its three axes and no further codimension-two contribution.
The good finite logarithmic caps have no residual critical
cycle after correcting the target.  Hence, with the complex
orientation and the spin normalization of D.2.44--D.2.45,
\[
                 \lambda(TX,\rho)=-\operatorname{PD}(C).
 \tag{D.WI6}
\]

Only a local covering is used to compute $[C]$.
In the positive period basis
$N b=a$, $N\gamma_1=\delta$, the weight-one cusp's literal
three-curve sum has regular-fibre preimage
$a\wedge b+b\wedge\delta+\delta\wedge\gamma_1$.
The weight-$d$ cusp is its degree-$d$ etale period cover, with
$\gamma_d=d\gamma_1$.  Transfer therefore gives the actual
preimage of its full $3d$-curve sum as
\[
 d\,a\wedge b+d\,b\wedge\delta+\delta\wedge\gamma_d.
 \tag{D.WI7}
\]
The angular vector $\delta$ here is fixed by
$N\gamma_d=d\delta$; this convention removes the possible
opposite choice for the original scalar-circle parameter.
Since $\xi_d$ has values $1,0,-6d$ on these three bivectors,
(D.WI7) has image $-5d\,s=0$.  Thus $C$ is integrally
nullhomologous and $c_2(TX)=\lambda(TX)=0$.
No covering of either finite logarithmic cap is assumed.

\emph{A finite-cap detector with its actual reference parity.}
If $d$ is odd, $2P_{\rm ref}=\lambda=0$ already gives
$P_{\rm ref}=0$.  Suppose $d$ is even.  In the order-four
finite cap, put $\beta=[a]$ and $\eta=[u_1]$ in its two
central base directions.  Its marked base lattice is
\[
 \langle\beta,\eta,g\mid m_1g=-2r_1\beta+n_1\eta\rangle .
 \tag{D.WI8}
\]
Choose integers $B,C$ with $m_1C=-2r_1+n_1B$ and define the
primitive character by $\chi(\beta)=1$, $\chi(\eta)=B$,
$\chi(g)=C$.  Such integers exist because $\gcd(m_1,n_1)=1$.
Its connected fibre $N$ in the central four-manifold is an
elliptic-torus mapping torus.  In its base kernel
$h=g-C\beta$ and $q=\eta-B\beta$ obey $m_1h=n_1q$.
A primitive generator has the root-loop exponent inverse to
$n_1$ modulo $m_1$, hence odd modulo four.  The elliptic
monodromy of $N$ is therefore the order-four matrix or its
inverse.  Its Wang sequence gives
\[
             H_1(N;\mathbb Z)=\mathbb Z\oplus\mathbb Z/2.
 \tag{D.WI9}
\]
Moreover $\chi(\delta)=-2$, so the actual scalar tangent
vector is transverse to $N$ in the reduced finite fibre.
Projecting $J\delta,\operatorname{Re}v_2,\operatorname{Im}v_2$
to $TN$ along $\delta$ identifies the stable tangent frame
induced by $\rho$ with an actual tangent parallelization.
The normal three-frame is $\delta$ followed by the two real
directions of $L$.

For an actual parallelization of a closed oriented
three-manifold, its stable framed coefficient modulo two is
$1+\dim H_1(N;\mathbb F_2)$.  The spin-compatible Kaplan
zero-/two-handle proof was given in the finite-II frame lemma:
the relative $\lambda$ reduces to the relative Euler number,
and the even intersection matrix has even mod-two rank.
It applies to every prescribed spin structure
\cite[Definition 1.7 and Theorem 3.1, pp.239,243--246]{mapsKaplan79}.
Equation (D.WI9) therefore gives reference parity
\[
                         \epsilon=[N]_\rho=1\pmod2 .
 \tag{D.WI10}
\]

Retain the actual Thom chains, not only their matrix classes.
Let $K_{a0},K_{a1}$ be the first columns of the two blocks in
(D.WE3).  Directly in its four target coordinates,
\[
                    d s=-2\ell P\,K_{a0}-kP\,K_{a1}.
 \tag{D.WI11}
\]
The chain at cap zero misses $N$, whereas the $a$-Thom
chain in cap one pairs with $N$ by $\chi(a)=1$.
Thus $\operatorname{lk}(N,s)=-kP/d$.
The determinant equation implies that $k$ is odd, so $N$
generates the two-primary summand of $H_3(X)$.
The specialization chains for (D.WI7) are supported in the
cusp and miss $N$.  Equations (D.WI7)--(D.WI11) then give
the actual integer link
\[
                          \ell(N,C)=5kP,
 \tag{D.WI12}
\]
up to the common harmless orientation sign.  It is odd.

Finally the relative refined-class computation extends with
reference parity $\epsilon$ in the form
\[
 \operatorname{lk}(N,\operatorname{PD}P_{\rm ref})
                 =-\frac{\ell(N,C)}2-\frac{\epsilon}2
                                      \quad\text{in }\mathbb Q/\mathbb Z .
 \tag{D.WI13}
\]
Indeed, choose Zhubr's admissible three-frame on a genuinely
embedded torsion sphere obtained from an odd multiple of $N$
by spin-compatible ambient two-handle surgery.  Its complement
spin three-frame extends over the trace.  The induced stable
framed class on $N$ is even.  The comparison with $\rho$ has
Spin(6) degree $q_0\equiv\epsilon\pmod2$.
For a bounding four-chain $B_4$ with boundary
$\operatorname{ord}(N)N$, the exact relative transgression is
$-B_4\cdot C=2q+\operatorname{ord}(N)q_0$.
Division of this integer identity gives (D.WI13), retaining
the actual rational link before reduction modulo one.
This is the proof of the relative lemma in D.2.45 with its
reference-parity term restored, not an application of the
unproved general change-of-frame conjecture.
Free $H_2$ causes no change: $C$ is nullhomologous, and
$2P_{\rm ref}=\lambda=0$ already kills the free part.

The two odd halves in (D.WI13) cancel by (D.WI10) and
(D.WI12).  The two-primary part of $P_{\rm ref}$ is zero,
and its odd-primary part was already zero from
$2P_{\rm ref}=0$.  This proves (D.WI2).
\end{proof}
